\documentclass[11pt]{article}
\usepackage[margin=1in]{geometry}
\usepackage{amsmath,amssymb,booktabs,microtype,bm}
\usepackage[hidelinks]{hyperref}
\usepackage{xcolor}
\newcommand{\Sym}{\mathrm{Sym}}
\newcommand{\tr}{\mathrm{tr}}
\newcommand{\Lie}{\mathcal{L}}
\newcommand{\norm}[1]{\lVert #1\rVert}
\newcommand{\code}[1]{\texttt{#1}}

\title{How the $N=2,3,4$ numerics were actually performed:\\
a full and honest methodology note}
\author{Numerical companion to the $\Sigma$ analytic notes}
\date{\today}

\begin{document}
\maketitle

\begin{abstract}
This note documents, in deliberately unvarnished detail, exactly how the
numerical study of the RG--invariant entropy--maximal variety $\Sigma$ was
carried out for $N=2,3,4$ flavours: what equations were solved, how the
residual and its Jacobian were built, which solver and acceptance criteria were
used, how the solution sets were post--processed into the dimensions and order
parameters quoted in \code{numerics\_summary.pdf}, and which independent checks
validate each step. It then reflects critically on \emph{why} the numerics
cleanly reproduces the analytic loci at $N=2$ and $N=3$ but reveals a genuinely
richer structure at $N=4$ --- separating the conclusions that are robust from
those that are soft, and stating plainly what sampling can and cannot establish.
All numbers cited are the ones in \code{metrics.json}, regenerated by
\code{make\_summary.py} from the saved solution sets
\code{n2\_clean.npz}, \code{n3\_clean.npz}, \code{n4\_clean.npz}
($k_{\max}=3$; $196/89/94$ solutions).
\end{abstract}

\section{What is being computed}
\label{sec:object}

The degrees of freedom are the components of a fully symmetric quartic coupling
$\lambda_{ijkl}=\lambda_{(ijkl)}$ on $\mathbb R^N$, i.e.\ a point of
$\Sym^4(\mathbb R^N)$, with
\[
n_{\mathrm{vars}}=\binom{N+3}{4}=5,\,15,\,35 \quad\text{for}\quad N=2,3,4.
\]
In code the independent components are indexed by sorted $4$--tuples
(\code{make\_keys}), and the full tensor is reconstructed by symmetrising over
the distinct permutations of each key (\code{inflate}/\code{deflate} in
\code{locus\_taylor.py}). We work throughout in the one--loop units
$\epsilon=\kappa=1$, in which the beta function is
\begin{equation}
\beta_{ijkl}=-\lambda_{ijkl}
 +\lambda_{ijmn}\lambda_{mnkl}
 +\lambda_{ikmn}\lambda_{mnjl}
 +\lambda_{ilmn}\lambda_{mnjk}.
\label{eq:beta}
\end{equation}
The one--leg ($1|N$) reduced density matrix and its spin--2 traceless part are
\begin{equation}
(\rho'_{1|N})_{ij}=\lambda_{iklm}\lambda_{jklm},\qquad
T_{ij}=(\rho'_{1|N})_{ij}-\tfrac1N(\tr\rho')\,\delta_{ij}.
\label{eq:T}
\end{equation}
Entropy of the one--leg cut is maximal precisely on $\mathcal M_{1|N}=\{T=0\}$.
The object of interest is not $\mathcal M_{1|N}$ itself but the largest
RG--flow--invariant subset of it, i.e.\ the points whose entire forward flow
line stays entropy--maximal. Writing the flow derivatives of $T$ as
$T^{(k)}=(\beta\!\cdot\!\partial_\lambda)^k T = k!\,[t^k]\,T(\lambda(t))$, this is
the \emph{Lie chain}
\begin{equation}
\Sigma=\bigl\{\lambda:\ T^{(0)}=T^{(1)}=T^{(2)}=\cdots=0\bigr\}.
\label{eq:Sigma}
\end{equation}
Imposing only $T^{(0)}=0$ (first--order tangency) is known to be too weak: the
three--flavour note (Remark~1) exhibits points that satisfy $T=T^{(1)}=0$ yet
whose flow leaves $\mathcal M_{1|N}$ at cubic order in RG time,
$\norm{T(t)}\sim t^3$. We therefore solve the \emph{truncated} chain
\begin{equation}
T^{(0)}=T^{(1)}=T^{(2)}=T^{(3)}=0\qquad(k_{\max}=3),
\label{eq:chain}
\end{equation}
and separately verify (Section~\ref{sec:validation}) that pushing to
$k_{\max}=4,5$ does not shrink the solution set further at sampled points. This
$k_{\max}=3$ is an \emph{empirical} stabilisation order, not a proven point of
ideal--stabilisation; that caveat is carried throughout.

\paragraph{Why a unit--norm pin is legitimate.}
On the variety $T=0$ the linear part of \eqref{eq:beta} acts on $T$ as a pure
rescaling: $\beta\!\cdot\!\partial_\lambda T = -2T + (\text{quadratic-in-}\lambda)\!\cdot\!\partial T$,
so on $\{T=0\}$ every higher condition $T^{(k\ge1)}$ reduces to a
\emph{homogeneous} polynomial in $\lambda$. Hence $\Sigma$ is a cone (a union of
rays through the origin), and slicing it with the unit sphere
$\norm{\lambda}=1$ captures each ray exactly once. Operationally we append
$\norm{x}^2-1$ as the final residual; it both removes the trivial root
$\lambda=0$ and renders the system square (favourable for the
Levenberg--Marquardt driver). The reported manifold dimensions are
\emph{cone} dimensions (the pin row is dropped before the Jacobian rank is
taken).

\section{How the residual and Jacobian are built}
\label{sec:residual}

Two independent implementations of $T^{(0..k_{\max})}(\lambda)$ exist; they
agree, which is the first line of defence against coding error.

\paragraph{(i) Symbolic build (\code{locus\_1to3\_Nflav.py}).}
Here $\rho$, $T$, $\beta$ and each Lie derivative
$\Lie_\beta T=\sum_\alpha\beta_\alpha\,\partial T/\partial\lambda_\alpha$ are
formed symbolically in \code{sympy}; the residual and an \emph{exact} Jacobian
are then \code{lambdify}'d to NumPy. This is transparent but does not scale:
the symbolic expansion of $(\Lie_\beta)^k T$ blows up combinatorially and is
already painful at $k=3$, so it was used mainly for $N=2,3$ cross--checks and
for the $O(N)$ orbit/PCA diagnostics it also houses (\code{classify\_keys},
\code{flattening\_matrix}).

\paragraph{(ii) Taylor--mode build (\code{locus\_taylor.py}, production).}
Rather than expanding $(\Lie_\beta)^k T$ symbolically, we integrate the flow as
a truncated power series. Each coupling is carried as a degree--$K$ series
$\lambda_\alpha(t)=\sum_{k=0}^{K}a_k t^k$ (the \code{Taylor} class: addition,
truncated convolution for products, and an einsum that contracts the tensor
axes coefficient--by--coefficient along the $t$--axis). Matching $t$--powers in
$\dot\lambda=\beta(\lambda)$ gives the recursion
\begin{equation}
(k+1)\,a_{k+1}=[t^k]\,\beta\bigl(\lambda(t)\bigr),\qquad k=0,\dots,K-1,
\end{equation}
implemented in \code{evolve\_lambda}. Feeding $\lambda(t)$ through
\eqref{eq:T} as a series yields $T(\lambda(t))$, and the chain values are read
off directly as $T^{(k)}=k!\,[t^k]T(\lambda(t))$. Each evaluation costs
$O(N^8)$ for the contractions in \eqref{eq:beta}--\eqref{eq:T} times $O(K^2)$
for the truncated products --- with \emph{no} symbolic blow--up. The residual
vector is the concatenation, over $k=0,\dots,k_{\max}$, of the upper--triangular
entries of $T^{(k)}$ \emph{minus} the redundant $(N\!-\!1,N\!-\!1)$ entry (fixed
by tracelessness), followed by the norm pin.

\paragraph{Complex--step Jacobian.}
The Jacobian fed to the solver is computed by complex--step differentiation:
column $i$ is $\operatorname{Im} R(x+i h\,e_i)/h$ with $h=10^{-30}$. Because the
residual is evaluated in a dtype--preserving way (the \code{Taylor}/\code{inflate}
machinery propagates \code{complex128} faithfully), this returns the derivative
to machine precision with \emph{no} subtractive cancellation --- strictly
better than a finite difference, and the reason the LM steps are clean even
near the rank--deficient solution manifold. (The docstring's mention of a
finite--difference Jacobian is stale; the code in \code{make\_residual} uses the
complex--step \code{cs\_jacobian}.)

\section{The solver and acceptance criteria}
\label{sec:solver}

Solutions are found by \emph{multistart} Levenberg--Marquardt
(\code{scipy.optimize.least\_squares}), in \code{run\_taylor.py}:

\begin{itemize}
\item \textbf{Starts.} Each start $x_0$ is drawn from a standard normal in
$\mathbb R^{n_{\mathrm{vars}}}$ and projected to the unit sphere. The RNG is
seeded (default \code{seed=0}) for reproducibility.
\item \textbf{Method.} \code{'lm'} (MINPACK) is used when the residual count
$\ge n_{\mathrm{vars}}$ (true here once the chain and the pin are stacked);
otherwise \code{'trf'}. LM is the right tool because the joint locus is a
\emph{positive--dimensional} variety, so the Jacobian at a solution is
rank--deficient; the damping term handles that kernel gracefully.
\item \textbf{Tolerances.} \code{xtol=ftol=gtol=1e-15}, \code{max\_nfev}
$=200$ (occasionally $400$). The tight inner tolerances let LM ride down to the
floor of the residual.
\item \textbf{Acceptance.} A start is \emph{accepted} if
$\max|\text{polynomial residual}|<\texttt{tol}=10^{-6}$ and the recovered norm
lies in $(0.5,2.0)$; the accepted point is then rescaled to
$\norm{x}=1$. The norm window rejects runaway/degenerate fits; the
$10^{-6}$ floor is the nominal membership threshold for $\Sigma$.
\item \textbf{Checkpointing.} Solutions are written to \code{.npz} every $25$
starts so a kill or timeout never loses progress.
\end{itemize}

The accepted, de--duplicated points are the clean sets
\code{n2/n3/n4\_clean.npz}. Acceptance is on $\norm{\cdot}_\infty<10^{-6}$, but
the \emph{realised} residuals are far tighter for the bulk (see
Section~\ref{sec:results}); the one exception is a handful of $N=3$ outliers
discussed below.

\section{Validation: the residual is independently trusted}
\label{sec:validation}

Because every downstream number rests on the residual being correct, it is
checked against machinery that shares no code path (\code{\_reconcile.py}):

\begin{description}
\item[A1 --- flow consistency.] The Taylor reconstruction
$T(t)=\sum_k T^{(k)}t^k/k!$ is compared, at several $t$, against an
\emph{independent} adaptive Runge--Kutta integration (\code{solve\_ivp},
DOP853, \code{rtol}$=10^{-12}$) of $\dot\lambda=\beta(\lambda)$ followed by a
direct evaluation of $T$. Agreement to $\lesssim10^{-7}$. This is the key
check: it validates the entire Taylor evaluator against a black--box ODE solve.
\item[A2/A3 --- $O(N)$ equivariance.] The symmetric ansatz point
($\lambda_{iiii}=u$, $\lambda_{iijj}=w$, else $0$) kills the chain to
$\sim10^{-9}$, and so does a random $O(3)$ rotation of it --- confirming that
$\beta$ and $T$ are correctly $O(N)$--equivariant (residual magnitude invariant
under rotation).
\item[B --- dimension reconciliation ($N=3$).] At a generic locus point the
Jacobian nullity is $5$; the three $\mathfrak{so}(3)$ rotation tangents lie in
that kernel (relative $\norm{J v}\sim10^{-6}$) and are rank $3$; hence
transverse moduli $=5-3=2$, exactly the note's two parameters $(u,w)$.
Imposing $k_{\max}=4,5$ leaves the nullity at $5$ --- the empirical evidence
that the chain has stabilised by $k_{\max}=3$.
\item[C --- the solutions are the note's locus.] On genuine multistart
solutions the note's algebra holds: the spin--2 trace part
$q_{ij}=\lambda_{ijkk}-\tfrac1N(\tr)\delta\approx0$, the harmonic part
$H=\lambda-sD$ is traceless, and its isotropy defect
$C_{ij}=H_{iklm}H_{jklm}-\tfrac1N(\tr)\delta\approx0$.
\end{description}

\section{Post--processing: how the reported quantities are defined}
\label{sec:postproc}

For each accepted solution (\code{make\_summary.py}):
\begin{itemize}
\item \textbf{Manifold (cone) dimension} $=n_{\mathrm{vars}}-\mathrm{rank}\,J$,
with $J$ the pin--free Jacobian and rank determined by a \emph{gap} heuristic
(\code{gap\_rank}): the largest base--10 gap in the singular--value spectrum
below a $10^{-3}\sigma_1$ floor. This is robust when there is a clean spectral
gap and \emph{soft} (can move by one) when there is not --- a caveat we flag
explicitly.
\item \textbf{Orbit dimension} $=\mathrm{rank}$ of the
$\mathfrak{so}(N)$ rotation tangents at the point (the infinitesimal $O(N)$
action on $\lambda$, deflated to component space). Generic points have full
$\dim O(N)=\binom N2 = 1,3,6$.
\item \textbf{Moduli} $=$ manifold dim $-$ orbit dim: the genuine physical
parameter count after quotienting the flavour rotations.
\item \textbf{Spin--2 defect} $\norm{q}$ and \textbf{isotropy defect}
(normalised $\norm{C}$): the two order parameters from the analytic proof.
\item \textbf{$\beta$ norm}: to flag fixed points ($\norm\beta<10^{-6}$).
\item \textbf{$\Sigma$--orbit spectrum match}: whether the eigenvalue
multiplicities of the $\Sym^2$ flattening matrix $M_{(ij)(kl)}=\lambda_{ijkl}$
match the symmetric line's $[1,N{-}1,\tfrac{N(N-1)}2]$. Because this spectrum
is an $O(N)$ invariant, it is a \emph{rotation--proof} test of whether a point
is in the $O(N)$ orbit of the symmetric ansatz.
\end{itemize}
A PCA of the centred solution cloud is also reported, but it measures the
\emph{linear span} of the sampled points, not a manifold dimension, and is
treated only as a coarse sanity check (it overshoots the local dimension
whenever the union of $O(N)$ orbits sweeps out extra linear directions).

\section{Results, case by case}
\label{sec:results}

\begin{center}\begin{tabular}{ccclccc}\toprule
$N$ & $n_{\mathrm{vars}}$ & \#sols & manifold dim (count) & median &
$\max|T^{(0\text{--}3)}|_\Sigma$ & $\Sigma$--orbit \\ \midrule
2 & 5  & 196 & 3\,(196)            & 3  & $0$                  & 196/196 \\
3 & 15 & 89  & 5\,(86), 10\,(3)    & 5  & $1.4\times10^{-12}$  & 0/89 \\
4 & 35 & 94  & 9\,(27), 11\,(18), 14\,(49) & 14 & $6.8\times10^{-13}$ & 0/94 \\
\bottomrule\end{tabular}\end{center}

\paragraph{$N=2$ ($5$ couplings, $196$ solutions).}
Residuals reach $1.4\times10^{-12}$. Every solution has manifold dim $3$,
orbit dim $1$ ($=\dim O(2)$), and moduli $2$; the spin--2 defect is
$\norm q\le10^{-15}$ and the isotropy defect vanishes for all $196$. Decoding
the components as $(a,b,c,d,e)$, the dedicated check (\code{report\_N2}) finds
$\max|a-e|$ and $\max|b+d|$ both at the $10^{-12}$ level: the numerics lands
\emph{exactly} on the analytic odd--block locus
$\mathcal E_-=\{(a,b,c,d,e)=(a,b,c,-b,a)\}$ of the two--flavour note (its
eq.~43). No fixed points are hit ($\norm\beta\in[1.23,12.71]$), as expected for
random multistart on a positive--dimensional variety.

\paragraph{$N=3$ ($15$ couplings, $89$ solutions).}
The dominant stratum ($86/89$) has manifold dim $5$, orbit dim $3$, moduli $2$,
with $q=0$ \emph{and} isotropic ($C=0$) on all $86$ --- precisely the analytic
cubic/tetrahedral structure
\[
\underbrace{5}_{\text{embedded cone}}
=\underbrace{2}_{\text{moduli }(u,w)}+\underbrace{3}_{\dim O(3)},
\]
reproducing the three--flavour theorem that $\Sigma$ is the $O(3)$ orbit of the
two--parameter plane $\lambda_{iiii}=u,\ \lambda_{iijj}=w$, with the spin--2
part \emph{forced} to vanish (Step~1) and a unique isotropic harmonic orbit
(Step~2). The $\Sigma$--orbit count is $0/89$ because the \emph{generic}
solution is the broader cubic harmonic, not the special $SO(3)$ (symmetric--line)
point --- consistent, not contradictory. The remaining $3/89$ are flagged as a
``dim $10$'' stratum with $q\neq0$ up to $6\times10^{-3}$ and the loosest
residuals in the whole study ($\le3.5\times10^{-7}$). Given the analytic proof
that $N=3$ has \emph{exactly one} orbit, the honest reading is that these three
are under--converged near--solutions, not a second component: at slightly
elevated residual the gap--rank heuristic loses its spectral gap and reports a
spuriously low rank (hence a spuriously high dimension). They are a numerical
artefact of finite tolerance, fully consistent with $N=3$ uniqueness.

\paragraph{$N=4$ ($35$ couplings, $94$ solutions).}
Residuals are tight ($\le6.1\times10^{-12}$) and again no fixed points are hit
($\norm\beta\in[1.50,12.32]$). But the locus is now \emph{reducible}: the
gap--rank manifold dimension takes three values,
\begin{center}\begin{tabular}{ccccccc}\toprule
dim & \#sols & orbit & moduli & $q{=}0$ & $\max\norm q$ & iso $C{=}0$\\ \midrule
9  & 27 & 6 & 3 & 5/27  & $0.341$ & 5/27 \\
11 & 18 & 6 & 5 & 0/18  & $0.479$ & 0/18 \\
14 & 49 & 6 & 8 & 49/49 & $0$     & 49/49 \\
\bottomrule\end{tabular}\end{center}
The decisive, threshold--independent finding is the last--but--two column: the
spin--2 part $q$ is \emph{not} forced to zero. $40/94$ solutions carry $q\neq0$,
up to $\norm q=0.479$ --- far above any tolerance, so this is real, not noise.
The dim--$14$ stratum ($49$ solutions) is the four--flavour analogue of the
$N=3$ family ($q=0$, isotropic); the dim--$9$ and dim--$11$ strata are
genuinely different, with large $q$ and non--isotropic harmonic parts. The
symmetric ansatz still solves the chain ($6.8\times10^{-13}$), but the
$\Sigma$--orbit match is $0/94$: the locus is strictly larger than the
symmetric line's orbit, and contains inequivalent harmonic--isotropy strata.

\section{Reflection: clean match at $N=2,3$, genuine break at $N=4$}
\label{sec:reflection}

\paragraph{Why $N=2,3$ match.}
At $N=2$ and $N=3$ the analytic loci are not merely sampled but \emph{proven
complete}: $\mathcal E_-$ at $N=2$, and at $N=3$ the exhaustive solve of the
five harmonic--isotropy quadratics $C_{ij}=0$ on the nine--dimensional space of
harmonic tensors, which yields a single $O(3)$ orbit (the cubic harmonic). The
numerics reproduces both at the level of \emph{independent} corroboration: the
order parameters the proof predicts to vanish do vanish on essentially every
sampled point ($q=0$ and $C=0$ on $196/196$ at $N=2$ and $86/86$ on the main
$N=3$ stratum), the dimension arithmetic $5=2+3$ comes out of the Jacobian
rank, and the flow--consistency check A1 ties the residual to a black--box ODE
integration. The agreement is strong precisely because the analytic side is a
theorem and the numerical side is a methodologically independent confirmation
of it.

\paragraph{Why $N=4$ is different.}
The ``discrepancy'' at $N=4$ is \emph{not} numerics disagreeing with analytics.
It is that the clean $N\le3$ paradigm --- $q$ forced to zero, a single isotropic
harmonic orbit --- \emph{fails}, and both the numerics and the four--flavour
analytic note agree that it fails. Analytically, the strictly analogous
harmonic computation at $N=4$ produces a \emph{second} orbit, the
hypertetrahedral $S_5\times\mathbb Z_2$ fixed point, distinguished from the
linear family $\Sigma^{\mathrm{lin}}_4$ by its $\Sym^2$--flattening eigenvalue
spectrum (an $O(4)$ invariant, hence not removable by any rotation): $\Sigma$ at
$N=4$ is \emph{strictly larger} than $\Sigma^{\mathrm{lin}}_4$. Numerically, the
two faces of this are (i) $q\neq0$ is now allowed --- the $N=3$ Step--1 theorem
has no $N=4$ counterpart --- and (ii) the solution set splits into a $q=0$
isotropic stratum plus inequivalent $q\neq0$ strata. The numerics independently
discovers the same qualitative break the analytic note predicts.

\paragraph{What is robust, and what is soft.}
Being honest about the strength of each $N=4$ claim:
\begin{itemize}
\item \textbf{Robust.} The existence of solutions with $q\neq0$ ($\norm q$ up to
$0.479$, orders of magnitude above tolerance) is threshold--independent and
decisively breaks the $N=3$ $q=0$ theorem at $N=4$. The existence of at least
one $q=0$/isotropic stratum (the dim--$14$ family) alongside genuinely distinct
$q\neq0$ strata is equally robust. The tight residuals ($\le6\times10^{-12}$)
confirm these are bona fide chain solutions, not near--misses.
\item \textbf{Soft.} The \emph{fine} stratification --- the exact dimensions
$9/11/14$ and their counts $27/18/49$ --- rests on the gap--rank heuristic,
which can shift a dimension by one on near--degenerate spectra, and on a sample
of only $94$ accepted points. The word ``component'' here means ``stratum by
local Jacobian rank,'' not a proven irreducible component.
\item \textbf{Sampling limits.} $N=4$ convergence is hard
($\sim30\text{--}40\%$ of starts succeed), so rare strata may be
under--sampled or missed entirely. Multistart establishes \emph{existence} of
points with given properties and \emph{estimates} local dimensions; it cannot
prove the stratification is exhaustive, nor prove maximality/completeness of
$\Sigma$. The analytic note supplies the existence and rotation--inequivalence
of the hypertetrahedral component; the numerics cannot, by itself, certify
there are no further pieces.
\end{itemize}

\paragraph{A note on fixed points.}
No solution in any case sits on an RG fixed point ($\norm\beta>1$ throughout).
This is expected: the catalogued fixed points (Gaussian, Decoupled, Cubic,
Heisenberg at $N=3$; additionally the hypertetrahedral point at $N=4$) lie on
$\Sigma$ but form a positive--codimension subset, which random multistart on a
continuous variety hits with probability zero. The analytic notes locate them
exactly; the numerics samples the generic flowing part of the same variety.
There is no tension here.

\section{Reproducibility and honest limitations}
\label{sec:repro}

Every number in \code{numerics\_summary.pdf} and in this note is computed by
\code{make\_summary.py} from the saved \code{.npz} sets and dumped to
\code{metrics.json}; nothing is hand--typed. To regenerate: run
\code{run\_taylor.py} with the desired $N$, $k_{\max}=3$, and a fixed seed to
produce the solution set, then \code{make\_summary.py} to recompute the
diagnostics. The principal limitations, restated plainly:
(i)~the chain is truncated at $k_{\max}=3$ and its stabilisation is checked only
at sampled points, not proved;
(ii)~``dimension'' is a gap--based local Jacobian rank, $\pm1$ on degenerate
spectra;
(iii)~``component/stratum'' is not a proven irreducible component;
(iv)~$N=4$ is under--sampled relative to its complexity, so absence of a
stratum in the data is not evidence of its absence in $\Sigma$;
(v)~multistart proves existence and estimates dimension, never completeness.
With those boundaries drawn, the numerics does what it is meant to do: it
independently confirms the $N=2,3$ theorems and independently exposes the
$N=4$ break that the four--flavour analytic note then explains.

\end{document}
